\begin{easyappendix}{Specyfikacja środowiska sprzętowo-programowego}

Zestawienie stacji roboczej oraz platformy testów mobilnych przedstawiono w tab.~\ref{tab:zal1_srodowisko}. Wartości pokrywają się z tab.~\ref{tab:sprzet} i tab.~\ref{tab:stos_programowy}; tutaj zebrano je w jednym miejscu.

\begin{table}[H]
\centering
\caption{Specyfikacja środowiska sprzętowo-programowego wykorzystanego do treningu i testów. Źródło: opracowanie własne.}
\label{tab:zal1_srodowisko}
\begin{tabular}{ |p{4.5cm}||p{8.5cm}| }
 \hline
 \multicolumn{2}{|c|}{\textbf{Stacja robocza (trening i ewaluacja GPU)}} \\
 \hline
 Procesor (CPU) & AMD Ryzen 5 3600 (6 rdzeni, 12 wątków, $3{,}60$\,GHz) \\
 Pamięć RAM & $32$\,GB DDR4 $3200$\,MHz \\
 Karta graficzna (GPU) & NVIDIA GeForce RTX 4070 ($12$\,GB VRAM CUDA) \\
 Pamięć masowa & $1$\,TB NVMe PCIe Gen4 SSD \\
 System operacyjny & Windows 11 / PyTorch $2.x$ CUDA $12.x$ \\
 \hline
 \multicolumn{2}{|c|}{\textbf{Środowisko testów mobilnych (inferencja krawędziowa)}} \\
 \hline
 Smartfon testowy & Samsung Galaxy S23 Ultra (Android 16) \\
 Mobilna stacja robocza & Lenovo ThinkPad (Intel Core i7, $16$\,GB RAM, Quadro T2000 $4$\,GB) \\
 Framework TFLite & TensorFlow Lite FP16 ($38{,}58$\,MB) / LiteRT \\
 \hline
\end{tabular}
\end{table}

\end{easyappendix}

\begin{easyappendix}{Wykaz 147 gatunków grzybów objętych systemem}

W tab.~\ref{tab:zal2_gatunki} przedstawiono kompletny wykaz $147$ gatunków grzybów objętych systemem detekcji i klasyfikacji, z uwzględnieniem nazwy łacińskiej, polskiej oraz kategorii toksykologicznej.

\begin{longtable}{|r|l|l|l|}
\caption{Wykaz $147$ gatunków objętych systemem wraz z kategorią toksykologiczną. Źródło: opracowanie własne.\label{tab:zal2_gatunki}} \\
\hline
\textbf{Lp.} & \textbf{Nazwa łacińska} & \textbf{Nazwa polska} & \textbf{Kategoria} \\
\hline
\endfirsthead
\hline
\textbf{Lp.} & \textbf{Nazwa łacińska} & \textbf{Nazwa polska} & \textbf{Kategoria} \\
\hline
\endhead
1 & \textit{Agaricus campestris} & Pieczarka łąkowa & Jadalny \\
2 & \textit{Amanita muscaria} & Muchomor czerwony & Trujący \\
3 & \textit{Amanita pantherina} & Muchomor plamisty & Śmiertelnie trujący \\
4 & \textit{Amanita phalloides} & Muchomor sromotnikowy & Śmiertelnie trujący \\
5 & \textit{Armillaria mellea} & Opieńka miodowa & Jadalny \\
6 & \textit{Auricularia auricula-judae} & Uszak bzowy & Jadalny \\
7 & \textit{Boletus edulis} & Borowik szlachetny & Jadalny \\
8 & \textit{Boletus reticulatus} & Borowik usiatkowany & Jadalny \\
9 & \textit{Calvatia gigantea} & Czasznica olbrzymia & Jadalny \\
10 & \textit{Cantharellus cibarius} & Pieprznik jadalny (kurka) & Jadalny \\
11 & \textit{Clavulina coralloides} & Gałęziak koralowaty & Jadalny \\
12 & \textit{Clitocybe nebularis} & Lejkówka szarawa & Niejadalny \\
13 & \textit{Coprinus comatus} & Czernidłak kołpakowaty & Jadalny \\
14 & \textit{Craterellus cornucopioides} & Lejkowiec dęty & Jadalny \\
15 & \textit{Daedaleopsis confragosa} & Gmatwica chropowata & Niejadalny \\
16 & \textit{Exidia glandulosa} & Kisielnica trzonkowa & Niejadalny \\
17 & \textit{Fistulina hepatica} & Ozorek dębowy & Jadalny \\
18 & \textit{Flammulina velutipes} & Płomiennica zimowa & Jadalny \\
19 & \textit{Fomes fomentarius} & Hubiak pospolity & Niejadalny \\
20 & \textit{Ganoderma lucidum} & Lakownica żółtawa & Niejadalny \\
21 & \textit{Gyromitra esculenta} & Piestrzenica kasztanowata & Śmiertelnie trujący \\
22 & \textit{Hydnum repandum} & Kolczak obłączasty & Jadalny \\
23 & \textit{Hygrophorus marzuolus} & Wodnica marcowa & Jadalny \\
24 & \textit{Imleria badia} & Podgrzybek brunatny & Jadalny \\
25 & \textit{Lactarius deliciosus} & Mleczaj rydz & Jadalny \\
26 & \textit{Leccinum scabrum} & Koźlarz babka & Jadalny \\
27 & \textit{Lepista nuda} & Gąsówka fioletowawa & Jadalny \\
28 & \textit{Macrolepiota procera} & Czubajka kania & Jadalny \\
29 & \textit{Morchella esculenta} & Smardz jadalny & Jadalny \\
30 & \textit{Panellus stipticus} & Łycznik ochrowy & Niejadalny \\
31 & \textit{Piptoporus betulinus} & Porek brzozowy & Niejadalny \\
32 & \textit{Pleurotus ostreatus} & Boczniak ostrygowaty & Jadalny \\
33 & \textit{Russula cyanoxantha} & Gołąbek zielonawofioletowy & Jadalny \\
34 & \textit{Russula emetica} & Gołąbek wymiotny & Trujący \\
35 & \textit{Russula virescens} & Gołąbek zielonawy & Jadalny \\
36 & \textit{Sarcoscypha coccinea} & Czarka szkarłatna & Jadalny \\
37 & \textit{Schizophyllum commune} & Rozszczepka pospolita & Niejadalny \\
38 & \textit{Suillus luteus} & Maślak zwyczajny & Jadalny \\
39 & \textit{Trametes versicolor} & Wrośniak różnobarwny & Niejadalny \\
40 & \textit{Tremella mesenterica} & Trzęsak pomarańczowożółty & Niejadalny \\
41 & \textit{Tricholoma equestre} & Gąska zielonka & Jadalny \\
42 & \textit{Xerocomus chrysenteron} & Suchogrzybek złotopory & Jadalny \\
43 & \textit{Cantharellus tubaeformis} & Pieprznik trąbkowy & Jadalny \\
44 & \textit{Collybia dryophila} & Pieniążek dębowy & Jadalny \\
45 & \textit{Gymnopilus junonius} & Łysak wspaniały & Trujący \\
46 & \textit{Hygrocybe conica} & Wilgotnica czerniejąca & Niejadalny \\
47 & \textit{Laccaria laccata} & Lakówka pospolita & Jadalny \\
48 & \textit{Mycena galericulata} & Grzybówka hełmiasta & Niejadalny \\
49 & \textit{Mycena pura} & Grzybówka fioletowawa & Trujący \\
50 & \textit{Hypholoma fasciculare} & Maślanka wiązkowa & Trujący \\
51 & \textit{Pholiota squarrosa} & Łuskwiak nastroszony & Niejadalny \\
52 & \textit{Stropharia aeruginosa} & Pierścieniak grynszpanowy & Niejadalny \\
53 & \textit{Psilocybe semilanceata} & Łysiczka lancetowata & Trujący \\
54 & \textit{Inocybe geophylla} & Strzępiak ziemistoblaszkowy & Trujący \\
55 & \textit{Cortinarius armillatus} & Zasłonak tasiemkowaty & Niejadalny \\
56 & \textit{Cortinarius violaceus} & Zasłonak fioletowy & Jadalny \\
57 & \textit{Paxillus involutus} & Krowiak podwinięty & Trujący \\
58 & \textit{Chroogomphus rutilus} & Klejek czerwonawy & Jadalny \\
59 & \textit{Gomphidius glutinosus} & Klejówka świerkowa & Jadalny \\
60 & \textit{Boletus pinophilus} & Borowik sosnowy & Jadalny \\
61 & \textit{Leccinum aurantiacum} & Koźlarz czerwony & Jadalny \\
62 & \textit{Tylopilus felleus} & Goryczak żółciowy & Niejadalny \\
63 & \textit{Scleroderma citrinum} & Tęgoskór cytrynowy & Trujący \\
64 & \textit{Lycoperdon perlatum} & Purchawka chropowata & Jadalny \\
65 & \textit{Bovista plumbea} & Czasznica ołowiana & Jadalny \\
66 & \textit{Cyathus striatus} & Kubek prążkowany & Niejadalny \\
67 & \textit{Crucibulum laeve} & Kubecznik pospolity & Niejadalny \\
68 & \textit{Helvella crispa} & Piestrzyca kędzierzawa & Niejadalny \\
69 & \textit{Verpa bohemica} & Naparstniczka czeska & Jadalny \\
70 & \textit{Amanita rubescens} & Muchomor czerwieniejący & Jadalny \\
71 & \textit{Amanita citrina} & Muchomor cytrynowy & Niejadalny \\
72 & \textit{Amanita fulva} & Muchomor rdzawy & Jadalny \\
73 & \textit{Russula nigricans} & Gołąbek czarniawy & Niejadalny \\
74 & \textit{Russula xerampelina} & Gołąbek śledziowy & Jadalny \\
75 & \textit{Russula paludosa} & Gołąbek błotny & Jadalny \\
76 & \textit{Lactarius torminosus} & Mleczaj wełnianka & Niejadalny \\
77 & \textit{Lactarius turpis} & Mleczaj ohydny & Niejadalny \\
78 & \textit{Lactarius piperatus} & Mleczaj biel & Niejadalny \\
79 & \textit{Lactarius rufus} & Mleczaj rudy & Niejadalny \\
80 & \textit{Suillus granulatus} & Maślak ziarnisty & Jadalny \\
81 & \textit{Suillus grevillei} & Maślak żółty & Jadalny \\
82 & \textit{Xerocomellus pruinatus} & Podgrzybek aksamitny & Jadalny \\
83 & \textit{Pseudoboletus parasiticus} & Podgrzybek pasożytniczy & Jadalny \\
84 & \textit{Gyroporus castaneus} & Piaskowiec kasztanowaty & Jadalny \\
85 & \textit{Clitopilus prunulus} & Sadoch jadalny & Jadalny \\
86 & \textit{Entoloma sinuatum} & Dzwonkówka trująca & Śmiertelnie trujący \\
87 & \textit{Lepiota procera} & Czubajka kania & Jadalny \\
88 & \textit{Chlorophyllum rhacodes} & Czubajka czerwieniejąca & Jadalny \\
89 & \textit{Agaricus arvensis} & Pieczarka zaroślowa & Jadalny \\
90 & \textit{Marasmius oreades} & Twardzioszek przydrożny & Jadalny \\
91 & \textit{Volvariella gloiocephala} & Pochwiak pochwiasty & Jadalny \\
92 & \textit{Pluteus cervinus} & Drobnołuszczak jeleni & Jadalny \\
93 & \textit{Tubaria furfuracea} & Trąbka otrębiasta & Niejadalny \\
94 & \textit{Galerina marginata} & Hełmówka jadowita & Śmiertelnie trujący \\
95 & \textit{Phallus impudicus} & Sromotnik bezwstydny & Niejadalny \\
96 & \textit{Hygrophoropsis aurantiaca} & Lisówka pomarańczowa & Niejadalny \\
97 & \textit{Amanita virosa} & Muchomor jadowity & Śmiertelnie trujący \\
98 & \textit{Agaricus xanthodermus} & Pieczarka karbolowa & Trujący \\
99 & \textit{Rubroboletus satanas} & Krwistoborowik szatański & Trujący \\
100 & \textit{Omphalotus olearius} & Kielichowiec pomarańczowy & Trujący \\
101 & \textit{Caloboletus calopus} & Gorzkoborowik żółtopory & Niejadalny \\
102 & \textit{Lepiota cristata} & Czubajeczka cuchnąca & Trujący \\
103 & \textit{Tricholoma sulphureum} & Gąska siarkowa & Trujący \\
104 & \textit{Cortinarius rubellus} & Zasłonak rudy & Śmiertelnie trujący \\
105 & \textit{Lactarius helvus} & Mleczaj torfowcowy & Trujący \\
106 & \textit{Amanita excelsa} & Muchomor twardawy & Jadalny \\
107 & \textit{Amanita gemmata} & Muchomor narcyzowy & Trujący \\
108 & \textit{Amanita porphyria} & Muchomor porfirowy & Niejadalny \\
109 & \textit{Calvatia utriformis} & Czasznica czaszkowata & Jadalny \\
110 & \textit{Cerioporus squamosus} & Żagwiak łuskowaty & Jadalny \\
111 & \textit{Clitocybe odora} & Lejkówka anyżkowa & Jadalny \\
112 & \textit{Coprinellus micaceus} & Czernidłaczek błyszczący & Jadalny \\
113 & \textit{Coprinopsis atramentaria} & Czernidłówka pospolita & Niejadalny \\
114 & \textit{Crepidotus variabilis} & Trzewiczek zmienny & Niejadalny \\
115 & \textit{Geastrum triplex} & Gwiazdosz potrójny & Niejadalny \\
116 & \textit{Grifola frondosa} & Żagwica listkowata & Jadalny \\
117 & \textit{Gymnopus dryophilus} & Pieniążek dębowy & Jadalny \\
118 & \textit{Gymnopus fusipes} & Pieniążek wrzecionowatotrzonowy & Niejadalny \\
119 & \textit{Gymnopus peronatus} & Pieniążek palący & Niejadalny \\
120 & \textit{Hebeloma crustuliniforme} & Włośnianka rosista & Trujący \\
121 & \textit{Hypholoma capnoides} & Maślanka łagodna & Jadalny \\
122 & \textit{Hypholoma lateritium} & Maślanka ceglasta & Niejadalny \\
123 & \textit{Kuehneromyces mutabilis} & Łuszczak zmienny & Jadalny \\
124 & \textit{Lactarius blennius} & Mleczaj śluzowaty & Niejadalny \\
125 & \textit{Lactarius quietus} & Mleczaj dębowy & Niejadalny \\
126 & \textit{Lactarius vellereus} & Mleczaj chrząstka & Niejadalny \\
127 & \textit{Laetiporus sulphureus} & Żółciak siarkowy & Jadalny \\
128 & \textit{Lentinula edodes} & Twardnik japoński (Shiitake) & Jadalny \\
129 & \textit{Lycoperdon pyriforme} & Purchawka gruszkowata & Jadalny \\
130 & \textit{Mycena epipterygia} & Grzybówka lepka & Niejadalny \\
131 & \textit{Mycena haematopus} & Grzybówka krwista & Niejadalny \\
132 & \textit{Neoboletus luridiformis} & Krasnoborowik ceglastopory & Jadalny \\
133 & \textit{Panaeolus foenisecii} & Kołpaczek szorstkozarodnikowy & Trujący \\
134 & \textit{Pleurotus pulmonarius} & Boczniak łyżkowaty & Jadalny \\
135 & \textit{Pluteus salicinus} & Drobnołuszczak wierzbowy & Niejadalny \\
136 & \textit{Russula foetens} & Gołąbek śmierdzący & Niejadalny \\
137 & \textit{Russula ochroleuca} & Gołąbek ochrowy & Niejadalny \\
138 & \textit{Russula vesca} & Gołąbek wyborny & Jadalny \\
139 & \textit{Scleroderma verrucosum} & Tęgoskór brodawkowaty & Trujący \\
140 & \textit{Stereum hirsutum} & Skórnik szorstki & Niejadalny \\
141 & \textit{Stropharia rugosoannulata} & Pierścieniak uprawny & Jadalny \\
142 & \textit{Suillus bovinus} & Maślak sitarz & Jadalny \\
143 & \textit{Suillus variegatus} & Maślak pstrokaty & Jadalny \\
144 & \textit{Trametes gibbosa} & Wrośniak garbaty & Niejadalny \\
145 & \textit{Tricholoma scalpturatum} & Gąska żółknąca & Jadalny \\
146 & \textit{Tricholoma terreum} & Gąska ziemista & Jadalny \\
147 & \textit{Xylaria polymorpha} & Próchnilec maczugowaty & Niejadalny \\
\hline
\end{longtable}

\end{easyappendix}

\begin{easyappendix}{Galeria próbek 147 gatunków grzybów}

Poniżej: po jednej reprezentatywnej fotografii dla każdego z $147$ gatunków objętych systemem. Obrazy pochodzą ze zbioru uczącego (iNaturalist, \texttt{quality\_grade=research}) po walidacji PIL.

\begin{figure}[htbp]
\centering
  \begin{subfigure}[b]{0.18\textwidth}
    \centering
    \includegraphics[width=\textwidth, height=2.5cm, keepaspectratio]{samples_147/sample_000_Agaricus_campestris.jpg}
    \caption*{\tiny 1. \textit{Agaricus campestris}}
  \end{subfigure}%
  \hfill
  \begin{subfigure}[b]{0.18\textwidth}
    \centering
    \includegraphics[width=\textwidth, height=2.5cm, keepaspectratio]{samples_147/sample_001_Amanita_muscaria.jpg}
    \caption*{\tiny 2. \textit{Amanita muscaria}}
  \end{subfigure}%
  \hfill
  \begin{subfigure}[b]{0.18\textwidth}
    \centering
    \includegraphics[width=\textwidth, height=2.5cm, keepaspectratio]{samples_147/sample_002_Amanita_pantherina.jpg}
    \caption*{\tiny 3. \textit{Amanita pantherina}}
  \end{subfigure}%
  \hfill
  \begin{subfigure}[b]{0.18\textwidth}
    \centering
    \includegraphics[width=\textwidth, height=2.5cm, keepaspectratio]{samples_147/sample_003_Amanita_phalloides.jpg}
    \caption*{\tiny 4. \textit{Amanita phalloides}}
  \end{subfigure}%
  \hfill
  \begin{subfigure}[b]{0.18\textwidth}
    \centering
    \includegraphics[width=\textwidth, height=2.5cm, keepaspectratio]{samples_147/sample_004_Armillaria_mellea.jpg}
    \caption*{\tiny 5. \textit{Armillaria mellea}}
  \end{subfigure}%
  \hfill\\
  \begin{subfigure}[b]{0.18\textwidth}
    \centering
    \includegraphics[width=\textwidth, height=2.5cm, keepaspectratio]{samples_147/sample_005_Auricularia_auricula-judae.jpg}
    \caption*{\tiny 6. \textit{Auricularia auricula-judae}}
  \end{subfigure}%
  \hfill
  \begin{subfigure}[b]{0.18\textwidth}
    \centering
    \includegraphics[width=\textwidth, height=2.5cm, keepaspectratio]{samples_147/sample_006_Boletus_edulis.jpg}
    \caption*{\tiny 7. \textit{Boletus edulis}}
  \end{subfigure}%
  \hfill
  \begin{subfigure}[b]{0.18\textwidth}
    \centering
    \includegraphics[width=\textwidth, height=2.5cm, keepaspectratio]{samples_147/sample_007_Boletus_reticulatus.jpg}
    \caption*{\tiny 8. \textit{Boletus reticulatus}}
  \end{subfigure}%
  \hfill
  \begin{subfigure}[b]{0.18\textwidth}
    \centering
    \includegraphics[width=\textwidth, height=2.5cm, keepaspectratio]{samples_147/sample_008_Calvatia_gigantea.jpg}
    \caption*{\tiny 9. \textit{Calvatia gigantea}}
  \end{subfigure}%
  \hfill
  \begin{subfigure}[b]{0.18\textwidth}
    \centering
    \includegraphics[width=\textwidth, height=2.5cm, keepaspectratio]{samples_147/sample_009_Cantharellus_cibarius.jpg}
    \caption*{\tiny 10. \textit{Cantharellus cibarius}}
  \end{subfigure}%
  \hfill\\
  \begin{subfigure}[b]{0.18\textwidth}
    \centering
    \includegraphics[width=\textwidth, height=2.5cm, keepaspectratio]{samples_147/sample_010_Clavulina_coralloides.jpg}
    \caption*{\tiny 11. \textit{Clavulina coralloides}}
  \end{subfigure}%
  \hfill
  \begin{subfigure}[b]{0.18\textwidth}
    \centering
    \includegraphics[width=\textwidth, height=2.5cm, keepaspectratio]{samples_147/sample_011_Clitocybe_nebularis.jpg}
    \caption*{\tiny 12. \textit{Clitocybe nebularis}}
  \end{subfigure}%
  \hfill
  \begin{subfigure}[b]{0.18\textwidth}
    \centering
    \includegraphics[width=\textwidth, height=2.5cm, keepaspectratio]{samples_147/sample_012_Coprinus_comatus.jpg}
    \caption*{\tiny 13. \textit{Coprinus comatus}}
  \end{subfigure}%
  \hfill
  \begin{subfigure}[b]{0.18\textwidth}
    \centering
    \includegraphics[width=\textwidth, height=2.5cm, keepaspectratio]{samples_147/sample_013_Craterellus_cornucopioides.jpg}
    \caption*{\tiny 14. \textit{Craterellus cornucopioides}}
  \end{subfigure}%
  \hfill
  \begin{subfigure}[b]{0.18\textwidth}
    \centering
    \includegraphics[width=\textwidth, height=2.5cm, keepaspectratio]{samples_147/sample_014_Daedaleopsis_confragosa.jpg}
    \caption*{\tiny 15. \textit{Daedaleopsis confragosa}}
  \end{subfigure}%
  \hfill\\
  \begin{subfigure}[b]{0.18\textwidth}
    \centering
    \includegraphics[width=\textwidth, height=2.5cm, keepaspectratio]{samples_147/sample_015_Exidia_glandulosa.jpg}
    \caption*{\tiny 16. \textit{Exidia glandulosa}}
  \end{subfigure}%
  \hfill
  \begin{subfigure}[b]{0.18\textwidth}
    \centering
    \includegraphics[width=\textwidth, height=2.5cm, keepaspectratio]{samples_147/sample_016_Fistulina_hepatica.jpg}
    \caption*{\tiny 17. \textit{Fistulina hepatica}}
  \end{subfigure}%
  \hfill
  \begin{subfigure}[b]{0.18\textwidth}
    \centering
    \includegraphics[width=\textwidth, height=2.5cm, keepaspectratio]{samples_147/sample_017_Flammulina_velutipes.jpg}
    \caption*{\tiny 18. \textit{Flammulina velutipes}}
  \end{subfigure}%
  \hfill
  \begin{subfigure}[b]{0.18\textwidth}
    \centering
    \includegraphics[width=\textwidth, height=2.5cm, keepaspectratio]{samples_147/sample_018_Fomes_fomentarius.jpg}
    \caption*{\tiny 19. \textit{Fomes fomentarius}}
  \end{subfigure}%
  \hfill
  \begin{subfigure}[b]{0.18\textwidth}
    \centering
    \includegraphics[width=\textwidth, height=2.5cm, keepaspectratio]{samples_147/sample_019_Ganoderma_lucidum.jpg}
    \caption*{\tiny 20. \textit{Ganoderma lucidum}}
  \end{subfigure}%
  \hfill\\
  \begin{subfigure}[b]{0.18\textwidth}
    \centering
    \includegraphics[width=\textwidth, height=2.5cm, keepaspectratio]{samples_147/sample_020_Gyromitra_esculenta.jpg}
    \caption*{\tiny 21. \textit{Gyromitra esculenta}}
  \end{subfigure}%
  \hfill
  \begin{subfigure}[b]{0.18\textwidth}
    \centering
    \includegraphics[width=\textwidth, height=2.5cm, keepaspectratio]{samples_147/sample_021_Hydnum_repandum.jpg}
    \caption*{\tiny 22. \textit{Hydnum repandum}}
  \end{subfigure}%
  \hfill
  \begin{subfigure}[b]{0.18\textwidth}
    \centering
    \includegraphics[width=\textwidth, height=2.5cm, keepaspectratio]{samples_147/sample_022_Hygrophorus_marzuolus.jpg}
    \caption*{\tiny 23. \textit{Hygrophorus marzuolus}}
  \end{subfigure}%
  \hfill
  \begin{subfigure}[b]{0.18\textwidth}
    \centering
    \includegraphics[width=\textwidth, height=2.5cm, keepaspectratio]{samples_147/sample_023_Imleria_badia.jpg}
    \caption*{\tiny 24. \textit{Imleria badia}}
  \end{subfigure}%
  \hfill
  \begin{subfigure}[b]{0.18\textwidth}
    \centering
    \includegraphics[width=\textwidth, height=2.5cm, keepaspectratio]{samples_147/sample_024_Lactarius_deliciosus.jpg}
    \caption*{\tiny 25. \textit{Lactarius deliciosus}}
  \end{subfigure}%
  \hfill\\
\end{figure}

\clearpage
\begin{figure}[htbp]
\centering
  \begin{subfigure}[b]{0.18\textwidth}
    \centering
    \includegraphics[width=\textwidth, height=2.5cm, keepaspectratio]{samples_147/sample_025_Leccinum_scabrum.jpg}
    \caption*{\tiny 26. \textit{Leccinum scabrum}}
  \end{subfigure}%
  \hfill
  \begin{subfigure}[b]{0.18\textwidth}
    \centering
    \includegraphics[width=\textwidth, height=2.5cm, keepaspectratio]{samples_147/sample_026_Lepista_nuda.jpg}
    \caption*{\tiny 27. \textit{Lepista nuda}}
  \end{subfigure}%
  \hfill
  \begin{subfigure}[b]{0.18\textwidth}
    \centering
    \includegraphics[width=\textwidth, height=2.5cm, keepaspectratio]{samples_147/sample_027_Macrolepiota_procera.jpg}
    \caption*{\tiny 28. \textit{Macrolepiota procera}}
  \end{subfigure}%
  \hfill
  \begin{subfigure}[b]{0.18\textwidth}
    \centering
    \includegraphics[width=\textwidth, height=2.5cm, keepaspectratio]{samples_147/sample_028_Morchella_esculenta.jpg}
    \caption*{\tiny 29. \textit{Morchella esculenta}}
  \end{subfigure}%
  \hfill
  \begin{subfigure}[b]{0.18\textwidth}
    \centering
    \includegraphics[width=\textwidth, height=2.5cm, keepaspectratio]{samples_147/sample_029_Panellus_stipticus.jpg}
    \caption*{\tiny 30. \textit{Panellus stipticus}}
  \end{subfigure}%
  \hfill\\
  \begin{subfigure}[b]{0.18\textwidth}
    \centering
    \includegraphics[width=\textwidth, height=2.5cm, keepaspectratio]{samples_147/sample_030_Piptoporus_betulinus.jpg}
    \caption*{\tiny 31. \textit{Piptoporus betulinus}}
  \end{subfigure}%
  \hfill
  \begin{subfigure}[b]{0.18\textwidth}
    \centering
    \includegraphics[width=\textwidth, height=2.5cm, keepaspectratio]{samples_147/sample_031_Pleurotus_ostreatus.jpg}
    \caption*{\tiny 32. \textit{Pleurotus ostreatus}}
  \end{subfigure}%
  \hfill
  \begin{subfigure}[b]{0.18\textwidth}
    \centering
    \includegraphics[width=\textwidth, height=2.5cm, keepaspectratio]{samples_147/sample_032_Russula_cyanoxantha.jpg}
    \caption*{\tiny 33. \textit{Russula cyanoxantha}}
  \end{subfigure}%
  \hfill
  \begin{subfigure}[b]{0.18\textwidth}
    \centering
    \includegraphics[width=\textwidth, height=2.5cm, keepaspectratio]{samples_147/sample_033_Russula_emetica.jpg}
    \caption*{\tiny 34. \textit{Russula emetica}}
  \end{subfigure}%
  \hfill
  \begin{subfigure}[b]{0.18\textwidth}
    \centering
    \includegraphics[width=\textwidth, height=2.5cm, keepaspectratio]{samples_147/sample_034_Russula_virescens.jpg}
    \caption*{\tiny 35. \textit{Russula virescens}}
  \end{subfigure}%
  \hfill\\
  \begin{subfigure}[b]{0.18\textwidth}
    \centering
    \includegraphics[width=\textwidth, height=2.5cm, keepaspectratio]{samples_147/sample_035_Sarcoscypha_coccinea.jpg}
    \caption*{\tiny 36. \textit{Sarcoscypha coccinea}}
  \end{subfigure}%
  \hfill
  \begin{subfigure}[b]{0.18\textwidth}
    \centering
    \includegraphics[width=\textwidth, height=2.5cm, keepaspectratio]{samples_147/sample_036_Schizophyllum_commune.jpg}
    \caption*{\tiny 37. \textit{Schizophyllum commune}}
  \end{subfigure}%
  \hfill
  \begin{subfigure}[b]{0.18\textwidth}
    \centering
    \includegraphics[width=\textwidth, height=2.5cm, keepaspectratio]{samples_147/sample_037_Suillus_luteus.jpg}
    \caption*{\tiny 38. \textit{Suillus luteus}}
  \end{subfigure}%
  \hfill
  \begin{subfigure}[b]{0.18\textwidth}
    \centering
    \includegraphics[width=\textwidth, height=2.5cm, keepaspectratio]{samples_147/sample_038_Trametes_versicolor.jpg}
    \caption*{\tiny 39. \textit{Trametes versicolor}}
  \end{subfigure}%
  \hfill
  \begin{subfigure}[b]{0.18\textwidth}
    \centering
    \includegraphics[width=\textwidth, height=2.5cm, keepaspectratio]{samples_147/sample_039_Tremella_mesenterica.jpg}
    \caption*{\tiny 40. \textit{Tremella mesenterica}}
  \end{subfigure}%
  \hfill\\
  \begin{subfigure}[b]{0.18\textwidth}
    \centering
    \includegraphics[width=\textwidth, height=2.5cm, keepaspectratio]{samples_147/sample_040_Tricholoma_equestre.jpg}
    \caption*{\tiny 41. \textit{Tricholoma equestre}}
  \end{subfigure}%
  \hfill
  \begin{subfigure}[b]{0.18\textwidth}
    \centering
    \includegraphics[width=\textwidth, height=2.5cm, keepaspectratio]{samples_147/sample_041_Xerocomus_chrysenteron.jpg}
    \caption*{\tiny 42. \textit{Xerocomus chrysenteron}}
  \end{subfigure}%
  \hfill
  \begin{subfigure}[b]{0.18\textwidth}
    \centering
    \includegraphics[width=\textwidth, height=2.5cm, keepaspectratio]{samples_147/sample_042_Cantharellus_tubaeformis.jpg}
    \caption*{\tiny 43. \textit{Cantharellus tubaeformis}}
  \end{subfigure}%
  \hfill
  \begin{subfigure}[b]{0.18\textwidth}
    \centering
    \includegraphics[width=\textwidth, height=2.5cm, keepaspectratio]{samples_147/sample_043_Collybia_dryophila.jpg}
    \caption*{\tiny 44. \textit{Collybia dryophila}}
  \end{subfigure}%
  \hfill
  \begin{subfigure}[b]{0.18\textwidth}
    \centering
    \includegraphics[width=\textwidth, height=2.5cm, keepaspectratio]{samples_147/sample_044_Gymnopilus_junonius.jpg}
    \caption*{\tiny 45. \textit{Gymnopilus junonius}}
  \end{subfigure}%
  \hfill\\
  \begin{subfigure}[b]{0.18\textwidth}
    \centering
    \includegraphics[width=\textwidth, height=2.5cm, keepaspectratio]{samples_147/sample_045_Hygrocybe_conica.jpg}
    \caption*{\tiny 46. \textit{Hygrocybe conica}}
  \end{subfigure}%
  \hfill
  \begin{subfigure}[b]{0.18\textwidth}
    \centering
    \includegraphics[width=\textwidth, height=2.5cm, keepaspectratio]{samples_147/sample_046_Laccaria_laccata.jpg}
    \caption*{\tiny 47. \textit{Laccaria laccata}}
  \end{subfigure}%
  \hfill
  \begin{subfigure}[b]{0.18\textwidth}
    \centering
    \includegraphics[width=\textwidth, height=2.5cm, keepaspectratio]{samples_147/sample_047_Mycena_galericulata.jpg}
    \caption*{\tiny 48. \textit{Mycena galericulata}}
  \end{subfigure}%
  \hfill
  \begin{subfigure}[b]{0.18\textwidth}
    \centering
    \includegraphics[width=\textwidth, height=2.5cm, keepaspectratio]{samples_147/sample_048_Mycena_pura.jpg}
    \caption*{\tiny 49. \textit{Mycena pura}}
  \end{subfigure}%
  \hfill
  \begin{subfigure}[b]{0.18\textwidth}
    \centering
    \includegraphics[width=\textwidth, height=2.5cm, keepaspectratio]{samples_147/sample_049_Hypholoma_fasciculare.jpg}
    \caption*{\tiny 50. \textit{Hypholoma fasciculare}}
  \end{subfigure}%
  \hfill\\
\end{figure}

\clearpage
\begin{figure}[htbp]
\centering
  \begin{subfigure}[b]{0.18\textwidth}
    \centering
    \includegraphics[width=\textwidth, height=2.5cm, keepaspectratio]{samples_147/sample_050_Pholiota_squarrosa.jpg}
    \caption*{\tiny 51. \textit{Pholiota squarrosa}}
  \end{subfigure}%
  \hfill
  \begin{subfigure}[b]{0.18\textwidth}
    \centering
    \includegraphics[width=\textwidth, height=2.5cm, keepaspectratio]{samples_147/sample_051_Stropharia_aeruginosa.jpg}
    \caption*{\tiny 52. \textit{Stropharia aeruginosa}}
  \end{subfigure}%
  \hfill
  \begin{subfigure}[b]{0.18\textwidth}
    \centering
    \includegraphics[width=\textwidth, height=2.5cm, keepaspectratio]{samples_147/sample_052_Psilocybe_semilanceata.jpg}
    \caption*{\tiny 53. \textit{Psilocybe semilanceata}}
  \end{subfigure}%
  \hfill
  \begin{subfigure}[b]{0.18\textwidth}
    \centering
    \includegraphics[width=\textwidth, height=2.5cm, keepaspectratio]{samples_147/sample_053_Inocybe_geophylla.jpg}
    \caption*{\tiny 54. \textit{Inocybe geophylla}}
  \end{subfigure}%
  \hfill
  \begin{subfigure}[b]{0.18\textwidth}
    \centering
    \includegraphics[width=\textwidth, height=2.5cm, keepaspectratio]{samples_147/sample_054_Cortinarius_armillatus.jpg}
    \caption*{\tiny 55. \textit{Cortinarius armillatus}}
  \end{subfigure}%
  \hfill\\
  \begin{subfigure}[b]{0.18\textwidth}
    \centering
    \includegraphics[width=\textwidth, height=2.5cm, keepaspectratio]{samples_147/sample_055_Cortinarius_violaceus.jpg}
    \caption*{\tiny 56. \textit{Cortinarius violaceus}}
  \end{subfigure}%
  \hfill
  \begin{subfigure}[b]{0.18\textwidth}
    \centering
    \includegraphics[width=\textwidth, height=2.5cm, keepaspectratio]{samples_147/sample_056_Paxillus_involutus.jpg}
    \caption*{\tiny 57. \textit{Paxillus involutus}}
  \end{subfigure}%
  \hfill
  \begin{subfigure}[b]{0.18\textwidth}
    \centering
    \includegraphics[width=\textwidth, height=2.5cm, keepaspectratio]{samples_147/sample_057_Chroogomphus_rutilus.jpg}
    \caption*{\tiny 58. \textit{Chroogomphus rutilus}}
  \end{subfigure}%
  \hfill
  \begin{subfigure}[b]{0.18\textwidth}
    \centering
    \includegraphics[width=\textwidth, height=2.5cm, keepaspectratio]{samples_147/sample_058_Gomphidius_glutinosus.jpg}
    \caption*{\tiny 59. \textit{Gomphidius glutinosus}}
  \end{subfigure}%
  \hfill
  \begin{subfigure}[b]{0.18\textwidth}
    \centering
    \includegraphics[width=\textwidth, height=2.5cm, keepaspectratio]{samples_147/sample_059_Boletus_pinophilus.jpg}
    \caption*{\tiny 60. \textit{Boletus pinophilus}}
  \end{subfigure}%
  \hfill\\
  \begin{subfigure}[b]{0.18\textwidth}
    \centering
    \includegraphics[width=\textwidth, height=2.5cm, keepaspectratio]{samples_147/sample_060_Leccinum_aurantiacum.jpg}
    \caption*{\tiny 61. \textit{Leccinum aurantiacum}}
  \end{subfigure}%
  \hfill
  \begin{subfigure}[b]{0.18\textwidth}
    \centering
    \includegraphics[width=\textwidth, height=2.5cm, keepaspectratio]{samples_147/sample_061_Tylopilus_felleus.jpg}
    \caption*{\tiny 62. \textit{Tylopilus felleus}}
  \end{subfigure}%
  \hfill
  \begin{subfigure}[b]{0.18\textwidth}
    \centering
    \includegraphics[width=\textwidth, height=2.5cm, keepaspectratio]{samples_147/sample_062_Scleroderma_citrinum.jpg}
    \caption*{\tiny 63. \textit{Scleroderma citrinum}}
  \end{subfigure}%
  \hfill
  \begin{subfigure}[b]{0.18\textwidth}
    \centering
    \includegraphics[width=\textwidth, height=2.5cm, keepaspectratio]{samples_147/sample_063_Lycoperdon_perlatum.jpg}
    \caption*{\tiny 64. \textit{Lycoperdon perlatum}}
  \end{subfigure}%
  \hfill
  \begin{subfigure}[b]{0.18\textwidth}
    \centering
    \includegraphics[width=\textwidth, height=2.5cm, keepaspectratio]{samples_147/sample_064_Bovista_plumbea.jpg}
    \caption*{\tiny 65. \textit{Bovista plumbea}}
  \end{subfigure}%
  \hfill\\
  \begin{subfigure}[b]{0.18\textwidth}
    \centering
    \includegraphics[width=\textwidth, height=2.5cm, keepaspectratio]{samples_147/sample_065_Cyathus_striatus.jpg}
    \caption*{\tiny 66. \textit{Cyathus striatus}}
  \end{subfigure}%
  \hfill
  \begin{subfigure}[b]{0.18\textwidth}
    \centering
    \includegraphics[width=\textwidth, height=2.5cm, keepaspectratio]{samples_147/sample_066_Crucibulum_laeve.jpg}
    \caption*{\tiny 67. \textit{Crucibulum laeve}}
  \end{subfigure}%
  \hfill
  \begin{subfigure}[b]{0.18\textwidth}
    \centering
    \includegraphics[width=\textwidth, height=2.5cm, keepaspectratio]{samples_147/sample_067_Helvella_crispa.jpg}
    \caption*{\tiny 68. \textit{Helvella crispa}}
  \end{subfigure}%
  \hfill
  \begin{subfigure}[b]{0.18\textwidth}
    \centering
    \includegraphics[width=\textwidth, height=2.5cm, keepaspectratio]{samples_147/sample_068_Verpa_bohemica.jpg}
    \caption*{\tiny 69. \textit{Verpa bohemica}}
  \end{subfigure}%
  \hfill
  \begin{subfigure}[b]{0.18\textwidth}
    \centering
    \includegraphics[width=\textwidth, height=2.5cm, keepaspectratio]{samples_147/sample_069_Amanita_rubescens.jpg}
    \caption*{\tiny 70. \textit{Amanita rubescens}}
  \end{subfigure}%
  \hfill\\
  \begin{subfigure}[b]{0.18\textwidth}
    \centering
    \includegraphics[width=\textwidth, height=2.5cm, keepaspectratio]{samples_147/sample_070_Amanita_citrina.jpg}
    \caption*{\tiny 71. \textit{Amanita citrina}}
  \end{subfigure}%
  \hfill
  \begin{subfigure}[b]{0.18\textwidth}
    \centering
    \includegraphics[width=\textwidth, height=2.5cm, keepaspectratio]{samples_147/sample_071_Amanita_fulva.jpg}
    \caption*{\tiny 72. \textit{Amanita fulva}}
  \end{subfigure}%
  \hfill
  \begin{subfigure}[b]{0.18\textwidth}
    \centering
    \includegraphics[width=\textwidth, height=2.5cm, keepaspectratio]{samples_147/sample_072_Russula_nigricans.jpg}
    \caption*{\tiny 73. \textit{Russula nigricans}}
  \end{subfigure}%
  \hfill
  \begin{subfigure}[b]{0.18\textwidth}
    \centering
    \includegraphics[width=\textwidth, height=2.5cm, keepaspectratio]{samples_147/sample_073_Russula_xerampelina.jpg}
    \caption*{\tiny 74. \textit{Russula xerampelina}}
  \end{subfigure}%
  \hfill
  \begin{subfigure}[b]{0.18\textwidth}
    \centering
    \includegraphics[width=\textwidth, height=2.5cm, keepaspectratio]{samples_147/sample_074_Russula_paludosa.jpg}
    \caption*{\tiny 75. \textit{Russula paludosa}}
  \end{subfigure}%
  \hfill\\
\end{figure}

\clearpage
\begin{figure}[htbp]
\centering
  \begin{subfigure}[b]{0.18\textwidth}
    \centering
    \includegraphics[width=\textwidth, height=2.5cm, keepaspectratio]{samples_147/sample_075_Lactarius_torminosus.jpg}
    \caption*{\tiny 76. \textit{Lactarius torminosus}}
  \end{subfigure}%
  \hfill
  \begin{subfigure}[b]{0.18\textwidth}
    \centering
    \includegraphics[width=\textwidth, height=2.5cm, keepaspectratio]{samples_147/sample_076_Lactarius_turpis.jpg}
    \caption*{\tiny 77. \textit{Lactarius turpis}}
  \end{subfigure}%
  \hfill
  \begin{subfigure}[b]{0.18\textwidth}
    \centering
    \includegraphics[width=\textwidth, height=2.5cm, keepaspectratio]{samples_147/sample_077_Lactarius_piperatus.jpg}
    \caption*{\tiny 78. \textit{Lactarius piperatus}}
  \end{subfigure}%
  \hfill
  \begin{subfigure}[b]{0.18\textwidth}
    \centering
    \includegraphics[width=\textwidth, height=2.5cm, keepaspectratio]{samples_147/sample_078_Lactarius_rufus.jpg}
    \caption*{\tiny 79. \textit{Lactarius rufus}}
  \end{subfigure}%
  \hfill
  \begin{subfigure}[b]{0.18\textwidth}
    \centering
    \includegraphics[width=\textwidth, height=2.5cm, keepaspectratio]{samples_147/sample_079_Suillus_granulatus.jpg}
    \caption*{\tiny 80. \textit{Suillus granulatus}}
  \end{subfigure}%
  \hfill\\
  \begin{subfigure}[b]{0.18\textwidth}
    \centering
    \includegraphics[width=\textwidth, height=2.5cm, keepaspectratio]{samples_147/sample_080_Suillus_grevillei.jpg}
    \caption*{\tiny 81. \textit{Suillus grevillei}}
  \end{subfigure}%
  \hfill
  \begin{subfigure}[b]{0.18\textwidth}
    \centering
    \includegraphics[width=\textwidth, height=2.5cm, keepaspectratio]{samples_147/sample_081_Xerocomellus_pruinatus.jpg}
    \caption*{\tiny 82. \textit{Xerocomellus pruinatus}}
  \end{subfigure}%
  \hfill
  \begin{subfigure}[b]{0.18\textwidth}
    \centering
    \includegraphics[width=\textwidth, height=2.5cm, keepaspectratio]{samples_147/sample_082_Pseudoboletus_parasiticus.jpg}
    \caption*{\tiny 83. \textit{Pseudoboletus parasiticus}}
  \end{subfigure}%
  \hfill
  \begin{subfigure}[b]{0.18\textwidth}
    \centering
    \includegraphics[width=\textwidth, height=2.5cm, keepaspectratio]{samples_147/sample_083_Gyroporus_castaneus.jpg}
    \caption*{\tiny 84. \textit{Gyroporus castaneus}}
  \end{subfigure}%
  \hfill
  \begin{subfigure}[b]{0.18\textwidth}
    \centering
    \includegraphics[width=\textwidth, height=2.5cm, keepaspectratio]{samples_147/sample_084_Clitopilus_prunulus.jpg}
    \caption*{\tiny 85. \textit{Clitopilus prunulus}}
  \end{subfigure}%
  \hfill\\
  \begin{subfigure}[b]{0.18\textwidth}
    \centering
    \includegraphics[width=\textwidth, height=2.5cm, keepaspectratio]{samples_147/sample_085_Entoloma_sinuatum.jpg}
    \caption*{\tiny 86. \textit{Entoloma sinuatum}}
  \end{subfigure}%
  \hfill
  \begin{subfigure}[b]{0.18\textwidth}
    \centering
    \includegraphics[width=\textwidth, height=2.5cm, keepaspectratio]{samples_147/sample_086_Lepiota_procera.jpg}
    \caption*{\tiny 87. \textit{Lepiota procera}}
  \end{subfigure}%
  \hfill
  \begin{subfigure}[b]{0.18\textwidth}
    \centering
    \includegraphics[width=\textwidth, height=2.5cm, keepaspectratio]{samples_147/sample_087_Chlorophyllum_rhacodes.jpg}
    \caption*{\tiny 88. \textit{Chlorophyllum rhacodes}}
  \end{subfigure}%
  \hfill
  \begin{subfigure}[b]{0.18\textwidth}
    \centering
    \includegraphics[width=\textwidth, height=2.5cm, keepaspectratio]{samples_147/sample_088_Agaricus_arvensis.jpg}
    \caption*{\tiny 89. \textit{Agaricus arvensis}}
  \end{subfigure}%
  \hfill
  \begin{subfigure}[b]{0.18\textwidth}
    \centering
    \includegraphics[width=\textwidth, height=2.5cm, keepaspectratio]{samples_147/sample_089_Marasmius_oreades.jpg}
    \caption*{\tiny 90. \textit{Marasmius oreades}}
  \end{subfigure}%
  \hfill\\
  \begin{subfigure}[b]{0.18\textwidth}
    \centering
    \includegraphics[width=\textwidth, height=2.5cm, keepaspectratio]{samples_147/sample_090_Volvariella_gloiocephala.jpg}
    \caption*{\tiny 91. \textit{Volvariella gloiocephala}}
  \end{subfigure}%
  \hfill
  \begin{subfigure}[b]{0.18\textwidth}
    \centering
    \includegraphics[width=\textwidth, height=2.5cm, keepaspectratio]{samples_147/sample_091_Pluteus_cervinus.jpg}
    \caption*{\tiny 92. \textit{Pluteus cervinus}}
  \end{subfigure}%
  \hfill
  \begin{subfigure}[b]{0.18\textwidth}
    \centering
    \includegraphics[width=\textwidth, height=2.5cm, keepaspectratio]{samples_147/sample_092_Tubaria_furfuracea.jpg}
    \caption*{\tiny 93. \textit{Tubaria furfuracea}}
  \end{subfigure}%
  \hfill
  \begin{subfigure}[b]{0.18\textwidth}
    \centering
    \includegraphics[width=\textwidth, height=2.5cm, keepaspectratio]{samples_147/sample_093_Galerina_marginata.jpg}
    \caption*{\tiny 94. \textit{Galerina marginata}}
  \end{subfigure}%
  \hfill
  \begin{subfigure}[b]{0.18\textwidth}
    \centering
    \includegraphics[width=\textwidth, height=2.5cm, keepaspectratio]{samples_147/sample_094_Phallus_impudicus.jpg}
    \caption*{\tiny 95. \textit{Phallus impudicus}}
  \end{subfigure}%
  \hfill\\
  \begin{subfigure}[b]{0.18\textwidth}
    \centering
    \includegraphics[width=\textwidth, height=2.5cm, keepaspectratio]{samples_147/sample_095_Hygrophoropsis_aurantiaca.jpg}
    \caption*{\tiny 96. \textit{Hygrophoropsis aurantiaca}}
  \end{subfigure}%
  \hfill
  \begin{subfigure}[b]{0.18\textwidth}
    \centering
    \includegraphics[width=\textwidth, height=2.5cm, keepaspectratio]{samples_147/sample_096_Amanita_virosa.jpg}
    \caption*{\tiny 97. \textit{Amanita virosa}}
  \end{subfigure}%
  \hfill
  \begin{subfigure}[b]{0.18\textwidth}
    \centering
    \includegraphics[width=\textwidth, height=2.5cm, keepaspectratio]{samples_147/sample_097_Agaricus_xanthodermus.jpg}
    \caption*{\tiny 98. \textit{Agaricus xanthodermus}}
  \end{subfigure}%
  \hfill
  \begin{subfigure}[b]{0.18\textwidth}
    \centering
    \includegraphics[width=\textwidth, height=2.5cm, keepaspectratio]{samples_147/sample_098_Rubroboletus_satanas.jpg}
    \caption*{\tiny 99. \textit{Rubroboletus satanas}}
  \end{subfigure}%
  \hfill
  \begin{subfigure}[b]{0.18\textwidth}
    \centering
    \includegraphics[width=\textwidth, height=2.5cm, keepaspectratio]{samples_147/sample_099_Omphalotus_olearius.jpg}
    \caption*{\tiny 100. \textit{Omphalotus olearius}}
  \end{subfigure}%
  \hfill\\
\end{figure}

\clearpage
\begin{figure}[htbp]
\centering
  \begin{subfigure}[b]{0.18\textwidth}
    \centering
    \includegraphics[width=\textwidth, height=2.5cm, keepaspectratio]{samples_147/sample_100_Caloboletus_calopus.jpg}
    \caption*{\tiny 101. \textit{Caloboletus calopus}}
  \end{subfigure}%
  \hfill
  \begin{subfigure}[b]{0.18\textwidth}
    \centering
    \includegraphics[width=\textwidth, height=2.5cm, keepaspectratio]{samples_147/sample_101_Lepiota_cristata.jpg}
    \caption*{\tiny 102. \textit{Lepiota cristata}}
  \end{subfigure}%
  \hfill
  \begin{subfigure}[b]{0.18\textwidth}
    \centering
    \includegraphics[width=\textwidth, height=2.5cm, keepaspectratio]{samples_147/sample_102_Tricholoma_sulphureum.jpg}
    \caption*{\tiny 103. \textit{Tricholoma sulphureum}}
  \end{subfigure}%
  \hfill
  \begin{subfigure}[b]{0.18\textwidth}
    \centering
    \includegraphics[width=\textwidth, height=2.5cm, keepaspectratio]{samples_147/sample_103_Cortinarius_rubellus.jpg}
    \caption*{\tiny 104. \textit{Cortinarius rubellus}}
  \end{subfigure}%
  \hfill
  \begin{subfigure}[b]{0.18\textwidth}
    \centering
    \includegraphics[width=\textwidth, height=2.5cm, keepaspectratio]{samples_147/sample_104_Lactarius_helvus.jpg}
    \caption*{\tiny 105. \textit{Lactarius helvus}}
  \end{subfigure}%
  \hfill\\
  \begin{subfigure}[b]{0.18\textwidth}
    \centering
    \includegraphics[width=\textwidth, height=2.5cm, keepaspectratio]{samples_147/sample_105_Amanita_excelsa.jpg}
    \caption*{\tiny 106. \textit{Amanita excelsa}}
  \end{subfigure}%
  \hfill
  \begin{subfigure}[b]{0.18\textwidth}
    \centering
    \includegraphics[width=\textwidth, height=2.5cm, keepaspectratio]{samples_147/sample_106_Amanita_gemmata.jpg}
    \caption*{\tiny 107. \textit{Amanita gemmata}}
  \end{subfigure}%
  \hfill
  \begin{subfigure}[b]{0.18\textwidth}
    \centering
    \includegraphics[width=\textwidth, height=2.5cm, keepaspectratio]{samples_147/sample_107_Amanita_porphyria.jpg}
    \caption*{\tiny 108. \textit{Amanita porphyria}}
  \end{subfigure}%
  \hfill
  \begin{subfigure}[b]{0.18\textwidth}
    \centering
    \includegraphics[width=\textwidth, height=2.5cm, keepaspectratio]{samples_147/sample_108_Calvatia_utriformis.jpg}
    \caption*{\tiny 109. \textit{Calvatia utriformis}}
  \end{subfigure}%
  \hfill
  \begin{subfigure}[b]{0.18\textwidth}
    \centering
    \includegraphics[width=\textwidth, height=2.5cm, keepaspectratio]{samples_147/sample_109_Cerioporus_squamosus.jpg}
    \caption*{\tiny 110. \textit{Cerioporus squamosus}}
  \end{subfigure}%
  \hfill\\
  \begin{subfigure}[b]{0.18\textwidth}
    \centering
    \includegraphics[width=\textwidth, height=2.5cm, keepaspectratio]{samples_147/sample_110_Clitocybe_odora.jpg}
    \caption*{\tiny 111. \textit{Clitocybe odora}}
  \end{subfigure}%
  \hfill
  \begin{subfigure}[b]{0.18\textwidth}
    \centering
    \includegraphics[width=\textwidth, height=2.5cm, keepaspectratio]{samples_147/sample_111_Coprinellus_micaceus.jpg}
    \caption*{\tiny 112. \textit{Coprinellus micaceus}}
  \end{subfigure}%
  \hfill
  \begin{subfigure}[b]{0.18\textwidth}
    \centering
    \includegraphics[width=\textwidth, height=2.5cm, keepaspectratio]{samples_147/sample_112_Coprinopsis_atramentaria.jpg}
    \caption*{\tiny 113. \textit{Coprinopsis atramentaria}}
  \end{subfigure}%
  \hfill
  \begin{subfigure}[b]{0.18\textwidth}
    \centering
    \includegraphics[width=\textwidth, height=2.5cm, keepaspectratio]{samples_147/sample_113_Crepidotus_variabilis.jpg}
    \caption*{\tiny 114. \textit{Crepidotus variabilis}}
  \end{subfigure}%
  \hfill
  \begin{subfigure}[b]{0.18\textwidth}
    \centering
    \includegraphics[width=\textwidth, height=2.5cm, keepaspectratio]{samples_147/sample_114_Geastrum_triplex.jpg}
    \caption*{\tiny 115. \textit{Geastrum triplex}}
  \end{subfigure}%
  \hfill\\
  \begin{subfigure}[b]{0.18\textwidth}
    \centering
    \includegraphics[width=\textwidth, height=2.5cm, keepaspectratio]{samples_147/sample_115_Grifola_frondosa.jpg}
    \caption*{\tiny 116. \textit{Grifola frondosa}}
  \end{subfigure}%
  \hfill
  \begin{subfigure}[b]{0.18\textwidth}
    \centering
    \includegraphics[width=\textwidth, height=2.5cm, keepaspectratio]{samples_147/sample_116_Gymnopus_dryophilus.jpg}
    \caption*{\tiny 117. \textit{Gymnopus dryophilus}}
  \end{subfigure}%
  \hfill
  \begin{subfigure}[b]{0.18\textwidth}
    \centering
    \includegraphics[width=\textwidth, height=2.5cm, keepaspectratio]{samples_147/sample_117_Gymnopus_fusipes.jpg}
    \caption*{\tiny 118. \textit{Gymnopus fusipes}}
  \end{subfigure}%
  \hfill
  \begin{subfigure}[b]{0.18\textwidth}
    \centering
    \includegraphics[width=\textwidth, height=2.5cm, keepaspectratio]{samples_147/sample_118_Gymnopus_peronatus.jpg}
    \caption*{\tiny 119. \textit{Gymnopus peronatus}}
  \end{subfigure}%
  \hfill
  \begin{subfigure}[b]{0.18\textwidth}
    \centering
    \includegraphics[width=\textwidth, height=2.5cm, keepaspectratio]{samples_147/sample_119_Hebeloma_crustuliniforme.jpg}
    \caption*{\tiny 120. \textit{Hebeloma crustuliniforme}}
  \end{subfigure}%
  \hfill\\
  \begin{subfigure}[b]{0.18\textwidth}
    \centering
    \includegraphics[width=\textwidth, height=2.5cm, keepaspectratio]{samples_147/sample_120_Hypholoma_capnoides.jpg}
    \caption*{\tiny 121. \textit{Hypholoma capnoides}}
  \end{subfigure}%
  \hfill
  \begin{subfigure}[b]{0.18\textwidth}
    \centering
    \includegraphics[width=\textwidth, height=2.5cm, keepaspectratio]{samples_147/sample_121_Hypholoma_lateritium.jpg}
    \caption*{\tiny 122. \textit{Hypholoma lateritium}}
  \end{subfigure}%
  \hfill
  \begin{subfigure}[b]{0.18\textwidth}
    \centering
    \includegraphics[width=\textwidth, height=2.5cm, keepaspectratio]{samples_147/sample_122_Kuehneromyces_mutabilis.jpg}
    \caption*{\tiny 123. \textit{Kuehneromyces mutabilis}}
  \end{subfigure}%
  \hfill
  \begin{subfigure}[b]{0.18\textwidth}
    \centering
    \includegraphics[width=\textwidth, height=2.5cm, keepaspectratio]{samples_147/sample_123_Lactarius_blennius.jpg}
    \caption*{\tiny 124. \textit{Lactarius blennius}}
  \end{subfigure}%
  \hfill
  \begin{subfigure}[b]{0.18\textwidth}
    \centering
    \includegraphics[width=\textwidth, height=2.5cm, keepaspectratio]{samples_147/sample_124_Lactarius_quietus.jpg}
    \caption*{\tiny 125. \textit{Lactarius quietus}}
  \end{subfigure}%
  \hfill\\
\end{figure}

\clearpage
\begin{figure}[htbp]
\centering
  \begin{subfigure}[b]{0.18\textwidth}
    \centering
    \includegraphics[width=\textwidth, height=2.5cm, keepaspectratio]{samples_147/sample_125_Lactarius_vellereus.jpg}
    \caption*{\tiny 126. \textit{Lactarius vellereus}}
  \end{subfigure}%
  \hfill
  \begin{subfigure}[b]{0.18\textwidth}
    \centering
    \includegraphics[width=\textwidth, height=2.5cm, keepaspectratio]{samples_147/sample_126_Laetiporus_sulphureus.jpg}
    \caption*{\tiny 127. \textit{Laetiporus sulphureus}}
  \end{subfigure}%
  \hfill
  \begin{subfigure}[b]{0.18\textwidth}
    \centering
    \includegraphics[width=\textwidth, height=2.5cm, keepaspectratio]{samples_147/sample_127_Lentinula_edodes.jpg}
    \caption*{\tiny 128. \textit{Lentinula edodes}}
  \end{subfigure}%
  \hfill
  \begin{subfigure}[b]{0.18\textwidth}
    \centering
    \includegraphics[width=\textwidth, height=2.5cm, keepaspectratio]{samples_147/sample_128_Lycoperdon_pyriforme.jpg}
    \caption*{\tiny 129. \textit{Lycoperdon pyriforme}}
  \end{subfigure}%
  \hfill
  \begin{subfigure}[b]{0.18\textwidth}
    \centering
    \includegraphics[width=\textwidth, height=2.5cm, keepaspectratio]{samples_147/sample_129_Mycena_epipterygia.jpg}
    \caption*{\tiny 130. \textit{Mycena epipterygia}}
  \end{subfigure}%
  \hfill\\
  \begin{subfigure}[b]{0.18\textwidth}
    \centering
    \includegraphics[width=\textwidth, height=2.5cm, keepaspectratio]{samples_147/sample_130_Mycena_haematopus.jpg}
    \caption*{\tiny 131. \textit{Mycena haematopus}}
  \end{subfigure}%
  \hfill
  \begin{subfigure}[b]{0.18\textwidth}
    \centering
    \includegraphics[width=\textwidth, height=2.5cm, keepaspectratio]{samples_147/sample_131_Neoboletus_luridiformis.jpg}
    \caption*{\tiny 132. \textit{Neoboletus luridiformis}}
  \end{subfigure}%
  \hfill
  \begin{subfigure}[b]{0.18\textwidth}
    \centering
    \includegraphics[width=\textwidth, height=2.5cm, keepaspectratio]{samples_147/sample_132_Panaeolus_foenisecii.jpg}
    \caption*{\tiny 133. \textit{Panaeolus foenisecii}}
  \end{subfigure}%
  \hfill
  \begin{subfigure}[b]{0.18\textwidth}
    \centering
    \includegraphics[width=\textwidth, height=2.5cm, keepaspectratio]{samples_147/sample_133_Pleurotus_pulmonarius.jpg}
    \caption*{\tiny 134. \textit{Pleurotus pulmonarius}}
  \end{subfigure}%
  \hfill
  \begin{subfigure}[b]{0.18\textwidth}
    \centering
    \includegraphics[width=\textwidth, height=2.5cm, keepaspectratio]{samples_147/sample_134_Pluteus_salicinus.jpg}
    \caption*{\tiny 135. \textit{Pluteus salicinus}}
  \end{subfigure}%
  \hfill\\
  \begin{subfigure}[b]{0.18\textwidth}
    \centering
    \includegraphics[width=\textwidth, height=2.5cm, keepaspectratio]{samples_147/sample_135_Russula_foetens.jpg}
    \caption*{\tiny 136. \textit{Russula foetens}}
  \end{subfigure}%
  \hfill
  \begin{subfigure}[b]{0.18\textwidth}
    \centering
    \includegraphics[width=\textwidth, height=2.5cm, keepaspectratio]{samples_147/sample_136_Russula_ochroleuca.jpg}
    \caption*{\tiny 137. \textit{Russula ochroleuca}}
  \end{subfigure}%
  \hfill
  \begin{subfigure}[b]{0.18\textwidth}
    \centering
    \includegraphics[width=\textwidth, height=2.5cm, keepaspectratio]{samples_147/sample_137_Russula_vesca.jpg}
    \caption*{\tiny 138. \textit{Russula vesca}}
  \end{subfigure}%
  \hfill
  \begin{subfigure}[b]{0.18\textwidth}
    \centering
    \includegraphics[width=\textwidth, height=2.5cm, keepaspectratio]{samples_147/sample_138_Scleroderma_verrucosum.jpg}
    \caption*{\tiny 139. \textit{Scleroderma verrucosum}}
  \end{subfigure}%
  \hfill
  \begin{subfigure}[b]{0.18\textwidth}
    \centering
    \includegraphics[width=\textwidth, height=2.5cm, keepaspectratio]{samples_147/sample_139_Stereum_hirsutum.jpg}
    \caption*{\tiny 140. \textit{Stereum hirsutum}}
  \end{subfigure}%
  \hfill\\
  \begin{subfigure}[b]{0.18\textwidth}
    \centering
    \includegraphics[width=\textwidth, height=2.5cm, keepaspectratio]{samples_147/sample_140_Stropharia_rugosoannulata.jpg}
    \caption*{\tiny 141. \textit{Stropharia rugosoannulata}}
  \end{subfigure}%
  \hfill
  \begin{subfigure}[b]{0.18\textwidth}
    \centering
    \includegraphics[width=\textwidth, height=2.5cm, keepaspectratio]{samples_147/sample_141_Suillus_bovinus.jpg}
    \caption*{\tiny 142. \textit{Suillus bovinus}}
  \end{subfigure}%
  \hfill
  \begin{subfigure}[b]{0.18\textwidth}
    \centering
    \includegraphics[width=\textwidth, height=2.5cm, keepaspectratio]{samples_147/sample_142_Suillus_variegatus.jpg}
    \caption*{\tiny 143. \textit{Suillus variegatus}}
  \end{subfigure}%
  \hfill
  \begin{subfigure}[b]{0.18\textwidth}
    \centering
    \includegraphics[width=\textwidth, height=2.5cm, keepaspectratio]{samples_147/sample_143_Trametes_gibbosa.jpg}
    \caption*{\tiny 144. \textit{Trametes gibbosa}}
  \end{subfigure}%
  \hfill
  \begin{subfigure}[b]{0.18\textwidth}
    \centering
    \includegraphics[width=\textwidth, height=2.5cm, keepaspectratio]{samples_147/sample_144_Tricholoma_scalpturatum.jpg}
    \caption*{\tiny 145. \textit{Tricholoma scalpturatum}}
  \end{subfigure}%
  \hfill\\
  \begin{subfigure}[b]{0.18\textwidth}
    \centering
    \includegraphics[width=\textwidth, height=2.5cm, keepaspectratio]{samples_147/sample_145_Tricholoma_terreum.jpg}
    \caption*{\tiny 146. \textit{Tricholoma terreum}}
  \end{subfigure}%
  \hfill
  \begin{subfigure}[b]{0.18\textwidth}
    \centering
    \includegraphics[width=\textwidth, height=2.5cm, keepaspectratio]{samples_147/sample_146_Xylaria_polymorpha.jpg}
    \caption*{\tiny 147. \textit{Xylaria polymorpha}}
  \end{subfigure}%
  \hfill
\end{figure}

\end{easyappendix}

